\introduction

На протяжении нескольких десятилетий рост производительности вычислительных
систем обеспечивался преимущественно за счёт увеличения тактовой частоты
процессоров. Однако в начале 2000-х годов данная тенденция столкнулась с
фундаментальными физическими ограничениями, и индустрия перешла к
многоядерным архитектурам~\cite{Sutter2005}. Сегодня флагманские серверные
процессоры содержат десятки и сотни вычислительных ядер, а параллелизм
превратился из узкоспециализированной темы в повседневную инженерную
реальность: веб-серверы, игровые движки, системы машинного обучения и базы
данных вынуждены явно проектироваться с расчётом на параллельное выполнение.

Ключевым инструментом организации параллельных вычислений в прикладных
системах является пул потоков~--- механизм, при котором фиксированный набор
заранее созданных потоков многократно используется для обработки поступающих
задач. Вместе с тем классические реализации пулов потоков с единой глобальной
очередью задач демонстрируют ограниченную масштабируемость: при увеличении
числа рабочих потоков единственная разделяемая очередь превращается в узкое
место, порождая конкуренцию и снижая эффективность. Стандартная библиотека
языка~C++ вплоть до версии~C++17 не предоставляет готовой реализации пула
потоков~\cite{Williams2019}, что вынуждает разработчиков либо прибегать к
крупным внешним зависимостям~--- таким как Intel Threading Building
Blocks,~--- либо реализовывать собственные решения.

Перспективным подходом к решению проблемы балансировки нагрузки является
механизм work-stealing (кражи задач), при котором каждый рабочий
поток владеет собственной локальной очередью, а незагруженный поток вправе
забирать задачи из очередей более загруженных коллег. Данный механизм обладает
строгим теоретическим обоснованием: Блюмофе и Лейзерсон доказали, что жадный
алгоритм work-stealing обеспечивает время выполнения $T_P \leq T_1 / P +
O(T_\infty)$, что является оптимальным с точностью до константного
множителя~\cite{Blumofe1999}. Несмотря на теоретическую зрелость механизма и
его практическое применение в промышленных библиотеках, в открытом доступе
отсутствуют реализации, одновременно удовлетворяющие следующим требованиям:
построены исключительно на стандарте C++11/14/17 без внешних зависимостей и
подкреплены систематическим экспериментальным исследованием характеристик
производительности. Изложенные обстоятельства определяют научную и
практическую актуальность настоящей работы.

Цель работы состоит в разработке библиотеки на языке~C++,
реализующей пул потоков с механизмом work-stealing, и в экспериментальном
исследовании её производительности в сравнении с классическим пулом потоков
при различных характеристиках нагрузки.

Достижение поставленной цели обеспечивается решением
следующих задач:
\begin{enumerate}
  \item Проектирование архитектуры библиотеки: разработать
    модульную архитектуру, включающую классический пул потоков с единой
    глобальной очередью и пул потоков с механизмом work-stealing, а также
    единый программный интерфейс, обеспечивающий взаимозаменяемость
    реализаций.

  \item Реализация двусторонней очереди: разработать структуру
    данных~--- двустороннюю очередь (deque)~--- для механизма work-stealing,
    обеспечивающую асимметричный доступ владельца и потоков-«воров» с
    корректной синхронизацией при конкурентном доступе.

  \item Реализация пулов потоков: построить обе реализации пула
    потоков исключительно на стандартных средствах языка~C++11/14/17 без
    внешних зависимостей в основной библиотеке.

  \item Верификация корректности: разработать автоматические тесты,
    охватывающие граничные случаи работы реализованных компонентов, и
    обеспечить отсутствие гонок данных посредством динамического анализа.

  \item Сравнительное исследование производительности: провести
    измерения производительности обеих реализаций для репрезентативных
    сценариев нагрузки при различном числе рабочих потоков и сформулировать
    практические рекомендации по применению.
\end{enumerate}

Объектом исследования являются пулы потоков как механизмы управления параллельным выполнением задач
в многоядерных вычислительных системах. Предметом исследования
выступают алгоритмические и реализационные характеристики механизма
work-stealing~--- в части его влияния на масштабируемость, пропускную
способность и потребление вычислительных ресурсов~--- в сравнении с
классической архитектурой с единой глобальной очередью.

Теоретические и методические основы работы. Теоретическую основу
исследования составляют следующие модели и результаты.

Модель многопоточных вычислений Блюмофе и
Лейзерсона~\cite{Blumofe1999}, представляющая параллельное вычисление в виде
ориентированного ациклического графа с параметрами работы~$T_1$ и критического
пути~$T_\infty$, используется для теоретического обоснования оптимальности
алгоритма work-stealing и формального определения целевых показателей
производительности.

Закон Амдала~\cite{Amdahl1967} применяется для анализа
принципиальных ограничений на масштабируемость параллельных реализаций и
обоснования требования к минимизации накладных расходов на синхронизацию.

Алгоритм динамической двусторонней очереди Чейза и
Лева~\cite{Chase2005} служит основой для проектирования структуры данных,
обеспечивающей асимметричный доступ владельца и воров.

Методические основы включают: практику параллельного программирования
на языке~C++ и применение примитивов \texttt{std::atomic},
\texttt{std::mutex}, \texttt{std::condition\_variable},
\texttt{std::future}, изложенную в книге Уильямса~\cite{Williams2019};
методы динамического анализа гонок данных посредством Thread~Sanitizer;
методику статистически корректного измерения производительности с использованием
библиотеки Google~Benchmark.

Настоящая работа вносит следующие результаты, соответствующие решённым задачам:

\begin{enumerate}
  \item Спроектирована модульная библиотека, предоставляющая единый
    шаблонный интерфейс \texttt{submit()} для обеих реализаций пула потоков,
    что обеспечивает корректность сравнительного исследования при одинаковых
    условиях.

  \item Реализована двусторонняя очередь \texttt{WorkStealingDeque}
    с асимметричным доступом: владелец работает по принципу LIFO с нижнего
    конца, потоки-воры~--- по принципу FIFO с верхнего конца, при
    корректной синхронизации посредством \texttt{std::mutex}.

  \item Реализованы оба пула потоков~--- \texttt{SimpleThreadPool}
    и \texttt{WorkStealingPool}~--- без внешних зависимостей, исключительно
    на средствах стандарта C++17.

  \item Верифицирована корректность всех компонентов посредством
    28~автоматических тестов с включённым Thread~Sanitizer; все тесты
    успешно пройдены без обнаруженных гонок данных. В ходе тестирования
    выявлены и устранены нетривиальные дефекты, в том числе гонка данных в
    деструкторе и переполнение беззнакового типа.

  \item Установлен фундаментальный компромисс между двумя
    реализациями: \texttt{SimpleThreadPool} обеспечивает меньшую латентность
    реакции на новые задачи (при 16~потоках в сценарии неравномерной
    нагрузки~--- 80{,}3~мс против 98{,}1~мс), тогда как
    \texttt{WorkStealingPool} потребляет в~3{,}8~раза меньше процессорного
    времени (2{,}75~мс против 10{,}4~мс), что определяет области
    предпочтительного применения каждой реализации.
\end{enumerate}

Научная новизна работы заключается в систематическом экспериментальном
исследовании влияния механизма ожидания задач на соотношение реального и
процессорного времени в реализациях пула потоков с work-stealing, а также в
выявлении условий, при которых блокирующая синхронизация в структуре
двусторонней очереди нивелирует теоретические преимущества механизма кражи
задач.

Разработанная библиотека снабжена воспроизводимой системой сборки на основе CMake с автоматическим
подключением зависимостей через механизм \texttt{FetchContent}. Библиотека
пригодна для непосредственного встраивания в проекты на языке~C++17 в качестве
замены \texttt{std::async} там, где требуется явное управление числом рабочих
потоков.